\input{../tex-inputs/latex-leseansicht-vorspann}

\section[Arthur Schnitzler und Olga Gussmann an Gustav Schwarzkopf, 12. 11. 1901]{L04449 Arthur Schnitzler und Olga Gussmann an Gustav Schwarzkopf, 12. 11. 1901}
\nopagebreak\mylabel{L04449v}
\rehead{ }\normalsize\beginnumbering\briefempfaengerindex{Schwarzkopf, Gustav@\textsc{Schwarzkopf, Gustav}!zzzSchnitzler, Olga@\emph{von Olga Schnitzler}!1901-11-121@{12. 11. 1901}|(be}\briefempfaengerindex{Schwarzkopf, Gustav@\textsc{Schwarzkopf, Gustav}!zzzSchnitzler, Arthur@\emph{von Arthur Schnitzler}!1901-11-121@{12. 11. 1901}|(be}
\toendnotes[C]{\smallbreak\pagebreak[2]}
\correspDesc{Versand  durch Arthur Schnitzler, Olga Gussmann am 12. 11. 1901 in Reichenau an der Rax
\newline{}Übermittlung  am 12. 11. 1901 in Reichenau an der Rax
\newline{}Zustellung  am 12. 11. 1901 in I., Innere Stadt
\newline{}Erhalt  durch Gustav Schwarzkopf am 12. 11. 1901 in Wien}\toendnotes[C]{\smallbreak}
\Standort{DLA, A:Schnitzler, HS.1985.1.1897.}
\physDesc{Bildpostkarte, 143 Zeichen
\newline{}Handschrift Arthur Schnitzler: Bleistift, deutsche Kurrent
\newline{}Handschrift Olga Schnitzler: Bleistift, lateinische Kurrent
\newline{}Versand: 1) Stempel: »\nobreak{}\oindex{Reichenau an der Rax@\textbf{Reichenau an der Rax}|pwk}\textcolor{gray}{Re}ichenau, 12 \textcolor{gray}{11} 01\nobreak{}«.   2) Stempel: »\nobreak{}\oindex{I., Innere Stadt@\textbf{I., Innere Stadt}|pwk}Wien 1/1 1, 12. 11. {[}01{]}, 6½–8 N, bestellt\nobreak{}«. }\pstart{}{\pb}Herrn Guſtav Schwarzkopf\pend{}\pstart{}Wien\oindex{Wien@\textbf{Wien}|pw}\pend{}\pstart{}I. Tiefer Graben 23\oindex{Wien@\textbf{Wien}!I., Innere Stadt@\textbf{I., Innere Stadt}!Tiefer Graben 23@\textbf{Tiefer Graben 23}|pw}.\pend{}{\bigskip}
\pstart
           \noindent{}{\pb}\textcolor{gray}{\textbf{Reichenau\oindex{Reichenau an der Rax@\textbf{Reichenau an der Rax}|pw}.}}\pend
           \vspace{1em}
\pstart
           {\pb}12. 11. 901\pend
           \vspace{0.5em}
\pstart
           Herrliches Wetter und herrliche Ruhe.\pend
           
\pstart
           Viele herzliche Grüße{\\[\baselineskip]}Ihr \spacefill\mbox{Arthur}\pend
           \leftskip=0em{}\selectlanguage{ngerman}\vspace{1em}
\pstart
           \noindent{}{[}hs. O. Schnitzler:{]} Viele Grüße!{\\}\spacefill\mbox{O. G.}\pend
           \selectlanguage{ngerman}\endnumbering\briefempfaengerindex{Schwarzkopf, Gustav@\textsc{Schwarzkopf, Gustav}!zzzSchnitzler, Olga@\emph{von Olga Schnitzler}!1901-11-121@{12. 11. 1901}|)be}\briefempfaengerindex{Schwarzkopf, Gustav@\textsc{Schwarzkopf, Gustav}!zzzSchnitzler, Arthur@\emph{von Arthur Schnitzler}!1901-11-121@{12. 11. 1901}|)be}\mylabel{L04449h}
\begin{anhang}
\end{anhang}\newcommand{\dateiname}{L04449}\newcommand{\titel}{Arthur Schnitzler und Olga Gussmann an Gustav Schwarzkopf, 12. 11. 1901}\newcommand{\editorInnen}{Herausgegeben von Jahnke, SelmaMüller, Martin Anton}\input{../tex-inputs/latex-leseansicht-abspann}
